% page 451 du fac-simile (image p-450 du scan)
% bloc 167.6 x 246.3 mm ; marge gauche 34.4 mm
\pgc{12.700mm}{0.000mm}{
\l{\cel{54}443}
\l{}
\l{\sou{placento}. I. (\sou{anat.}) Organo karnoza, sponjatra, qua adheras a}
\l{\cel{3}la utero, komunikante per la kordono umbilikala kun la feto,}
\l{\cel{3}per la vaskuli di qua la sango patrinala arivas a la feto e}
\l{\cel{3}retrovenas a la patrino. - II. (\sou{bot.}) Parto di la surfaco di}
\l{\cel{3}karpelo, sur qua esas insertita la ovuli.--DEFIS.}
\l{}
\l{\sou{plado}. Vazo di qua la fundo esas plata,quan onu pozas avan}
\l{\cel{3}singla repastonto o festinonto por recevar la nutrivi. - EFIS.}
\l{}
\l{\sou{plafono}. (\sou{arkitekt}.) Surfaco ye qua konsistas la parto supra di}
\l{\cel{3}chambro, salono, galerio, kirko, e c., super la kapi di}
\l{\cel{3}la personi qui esas en lu. - DFR.}
\l{}
\l{\sou{plago}. I. (\sou{fiziol}.)Laceruro, extera od interna, kun desorganiz-}
\l{\cel{3}eso di la tisui en parto di la korpo, produktita da vunduro}
\l{\cel{3}o kauzo morbifera. - II. (\sou{metaf}.)I. Persono o kozo konsiderata}
\l{\cel{3}kom instrumento per qua Deo frapas o kastigas la homi. II.}
\l{\cel{3}Kalamitato qua frapas populo. III. To quo esas funesta a per-}
\l{\cel{3}sono od a kozo. - DEFS.}
\l{}
\l{\sou{plajo}. Bendo, plu o min larja, de sulo alonge maro od oceano,}
\l{\cel{3}quan ta maro o ta oceano alterne kovras e deskovras. (Existas}
\l{\cel{3}plaji de sablo e plaji de rul-stoni). - FIS.}
\l{}
\l{\sou{plajiar}. (\sou{trans.}) (\sou{literaturo}).Furtar ek la verki da kunhomo.}
\l{\cel{3}- DEFIRS.}
\l{}
\l{\sou{plako}. Lamo de metalo - maxim-multa-kaze destinita aplikesor}
\l{\cel{3}adsur surfaco plana. EFIR.}
\l{}
\l{\sou{plano}. I. Surfaco sen kurviguro nek ondaro. - II. (\sou{geom.})}
\l{\cel{3}\sou{Surfaco plana}\cel{1}: sur qua rekto esas aplikebla omna-sinse. -}
\l{\cel{3}\sou{Angulo plana}\cel{1}: intersekuro da du rekti sama-plana. Figuro}
\l{\cel{3}\sou{plana}\cel{1}:\cel{2}trasita sur surfaco plana. - III. Desegnuro qua repre-}
\l{\cel{3}zentas la dispozeso di la kompozanti di urbo, di regiono, di}
\l{\cel{3}edifico, kun indiko di la proporcioni e di la plaso relativa di}
\l{\cel{3}la parti diversa. - DEFIRS.}
\l{}
\l{\sou{plando}. (\sou{anat.}) Parto di la pedo qua tushas sulo kande onu stacas}
\l{\cel{3}o kande onu marchas. - FIS.}
\l{}
\l{\sou{planeto}. I. (\sou{astron. antiqua})\cel{2}Astro vaganta (kontraste kun la}
\l{\cel{3}steli fixa). (\sou{astrol}.)Astro a qua onu imputas efiko a la}
\l{\cel{3}fato di la homo. - II. (\sou{astron. moderna}). Astro qua parkuras}
\l{\cel{3}la spaco cirkum Suno (situita ad un ek la foki). - DEFIRS.}
\l{}
\l{\sou{planimetro.}\cel{2}Instrumento per qua onu povas evaluar la mezuri di}
\l{\cel{3}areo plana. - DEFIRS.}
\l{}
\l{planimetrio. Parto di geometrio : la arto mezurar la surfaci}
\l{\cel{3}plana. - DEFIRS.}
\l{}
\l{planisfero. -(\sou{geogr}.)Mapo en qua la Ter-sfero dividesas en du}
\l{\cel{3}diametri, reprezentata da surfaco plana, egardinte la redukto}
\l{\cel{3}demandata da perspektivo. - DEFIRS.}
}
